%File: main.tex
%AAAI 2026 style final project report
%
% USAGE: Drop this file into the AAAI 2026 Author Kit directory
% (alongside aaai2026.sty and aaai2026.bst), or upload to Overleaf using
% the official AAAI 2026 template and replace the main .tex file with this one.
% Then: pdflatex main; bibtex main; pdflatex main; pdflatex main.

\documentclass[letterpaper]{article}
\usepackage{aaai2026}  % AAAI 2026 style
\usepackage{times}
\usepackage{helvet}
\usepackage{courier}
\usepackage[hyphens]{url}
\usepackage{graphicx}
\urlstyle{rm}
\def\UrlFont{\rm}
\usepackage{natbib}
\usepackage{caption}
\DeclareCaptionStyle{ruled}{labelfont=normalfont,labelsep=colon,strut=off}
\frenchspacing
\setlength{\pdfpagewidth}{8.5in}
\setlength{\pdfpageheight}{11in}

% Packages beyond what's in the AAAI style
\usepackage{amsmath}
\usepackage{amssymb}
\usepackage{booktabs}
\usepackage{algorithm}
\usepackage{algpseudocode}
\usepackage[section]{placeins}

% --- Float placement tuning ---
% AAAI's narrow two-column layout makes LaTeX very sensitive to float
% placement. These settings keep floats near their reference text and
% prevent the ``tables crammed against the bottom margin'' pattern.
\renewcommand{\topfraction}{0.9}
\renewcommand{\bottomfraction}{0.8}
\renewcommand{\textfraction}{0.07}
\renewcommand{\floatpagefraction}{0.75}
\setlength{\textfloatsep}{10pt plus 2pt minus 2pt}
\setlength{\intextsep}{10pt plus 2pt minus 2pt}
\renewcommand{\arraystretch}{1.08}

\pdfinfo{
/Title (Model-Based vs. Model-Free Reinforcement Learning in Partially Observable MiniGrid Environments)
/Author (George Wang, Shreyas Suresh Manoti)
/TemplateVersion (2026.1)
}

\setcounter{secnumdepth}{2}

\title{Model-Based vs.\ Model-Free Reinforcement Learning in\\
Partially Observable MiniGrid Environments}

\author{
    George Wang,\ Shreyas Suresh Manoti\\
}
\affiliations{
    Northeastern University\\
    CS 5180: Reinforcement Learning, Spring 2026\\
}

\begin{document}
\maketitle

\begin{abstract}
We compare a recurrent model-free baseline (Recurrent PPO) against a
deterministic world-model agent with Dyna-style imagination on three
partially observable MiniGrid tasks: \textsc{DoorKey-5x5} and
\textsc{DoorKey-8x8} (sequential object interaction at two difficulty
scales) and \textsc{MemoryS11} (long-term recall). Both agents share the
same CNN encoder and GRU belief backbone; the world-model agent adds a
latent-transition head, an observation-reconstruction head, optional
curiosity, and an imagination rollout that performs a second policy update
on simulated trajectories. After controlled experiments we find that both
methods solve \textsc{DoorKey-5x5} reliably (final mean return $\approx
0.96$ for both), though Recurrent PPO actually reaches the $0.95$
threshold about twice as fast as the world-model agent on this smallest
variant; on \textsc{MemoryS11} both methods plateau near mean return
$0.47$, which on this task corresponds to policies that do not reliably
use the remembered cue; and neither method solves \textsc{DoorKey-8x8}
reliably within the training budget ($0/2$ seeds reach goal-level return
under either method). The world-model additions produce one measurable
qualitative effect---a smoother, lower-variance learning curve on
\textsc{MemoryS11}---but they do not raise the return ceiling on any of
the three tasks. Our main contribution is a set of three
diagnosed failure modes in the model-based pipeline---imagination through
the transition model, insufficient entropy regularization, and detaching
the belief state from the policy gradient---each of which we show to be a
design trade-off rather than an implementation bug. A fourth contribution
is a \emph{temporal retention probe} on \textsc{MemoryS11} that decodes the
start-of-episode cue from the frozen belief state of a trained checkpoint
and finds that cue identity is not linearly recoverable past the cue room,
which mechanistically explains the $\approx 0.47$ plateau. Together these results
argue that auxiliary dynamics losses supplement, but cannot replace,
end-to-end gradient flow from the policy into the recurrent belief state,
and that value-only imagination without a learned reward model is a
brittle simplification of the Dreamer recipe.
\end{abstract}

\section{Introduction}

Partially observable environments are a long-standing challenge for
reinforcement learning (RL): the agent must infer latent state from a history
of observations rather than a single frame, and sparse rewards make credit
assignment across long horizons difficult. Two broad strategies dominate the
literature. \emph{Model-free recurrent policies} such as Recurrent PPO
\cite{schulman2017ppo} summarize history implicitly in a recurrent hidden
state that is shaped only by the task reward. \emph{Model-based} agents such
as PlaNet \cite{hafner2019planet} and Dreamer \cite{hafner2020dreamer} learn
an explicit dynamics model and use it either for planning or for generating
imagined rollouts on which the policy is further trained.

The promise of the model-based approach is two-fold. First, the dynamics
prediction objective acts as a representation-learning signal: the recurrent
state is pressured to capture structure that is useful for predicting the
future, not just for the next action. Second, imagined rollouts decouple
policy improvement from environment interaction, potentially improving sample
efficiency. Both claims come with caveats, which we examine empirically in
this work.

We study these two strategies on the \textsc{MiniGrid} suite
\cite{minigrid}, specifically \textsc{DoorKey-5x5-v0},
\textsc{DoorKey-8x8-v0}, and \textsc{MemoryS11-v0}. These tasks provide a
controlled setting in which partial observability arises in different ways:
\textsc{DoorKey} requires sequential object interaction (locate the key,
toggle the door, reach the goal), while \textsc{MemoryS11} requires
remembering an object observed at episode start and matching it with a cue
observed later. Including both sizes of \textsc{DoorKey} lets us isolate
the effect of exploration horizon from the effect of task structure: the
5$\times$5 and 8$\times$8 variants share the same observation space and
action space and differ only in room size and therefore in the number of
random actions required to stumble on the key-then-door interaction. All
three environments have sparse reward ($+1$ at the goal, $0$ otherwise,
minus a small per-step penalty) and 7$\times$7$\times$3 local egocentric
observations.

Our contributions are:

\begin{itemize}
    \item A controlled comparison of Recurrent PPO and a deterministic
    world-model agent (GRU belief state, latent-transition head,
    observation-reconstruction head, optional curiosity, Dyna-style
    imagination) under matched architecture and compute on
    \textsc{DoorKey-5x5}, \textsc{DoorKey-8x8}, and \textsc{MemoryS11}.
    \item Three diagnosed negative results---BPTT-through-transition
    imagination, under-regularized entropy on sparse-reward tasks, and
    detachment of the belief state from the policy gradient---each
    quantified with the training scalars that identified them and mapped to
    a concrete design implication.
    \item An empirical argument that a purely deterministic world model with
    value-only imagination is not a viable simplification of Dreamer: without
    a learned reward model, imagination gradients are high-variance and
    interfere destructively with the clean one-step prediction loss.
\end{itemize}

\section{Background}

\subsection{POMDP Formulation}

We model the environment as a partially observable MDP
$(\mathcal{S}, \mathcal{A}, \mathcal{O}, T, R, \Omega, \gamma)$ with hidden
state space $\mathcal{S}$ (agent pose, door/key status, goal location),
discrete action space $\mathcal{A}$ (turn left/right, move forward, pick up,
toggle), observation space $\mathcal{O}$ of 7$\times$7$\times$3 egocentric
RGB views, transition kernel $T(s' \mid s, a)$, sparse reward $R(s,a)$,
observation kernel $\Omega(o \mid s', a)$, and discount $\gamma \in [0,1)$.
The agent maintains a history $h_t = (o_1, a_1, \ldots, o_t)$ and seeks
a policy $\pi(a \mid h_t)$ maximizing
$\mathbb{E}\left[\sum_{t=0}^{\infty} \gamma^t R(s_t, a_t)\right]$.

\subsection{PPO and Recurrent PPO}

Proximal Policy Optimization \cite{schulman2017ppo} optimizes a clipped
surrogate objective
\begin{equation*}
    \mathcal{L}_{\text{PPO}} = \mathbb{E}_t\!\left[\min\!\big(r_t(\theta) \hat A_t,\;
    \text{clip}(r_t(\theta), 1{-}\epsilon, 1{+}\epsilon)\hat A_t\big)\right]
\end{equation*}
with importance ratio $r_t(\theta) = \pi_\theta(a_t \mid s_t) / \pi_{\theta_\text{old}}(a_t \mid s_t)$
and GAE advantages $\hat A_t$. Recurrent PPO additionally maintains a GRU
hidden state, detaches the state at episode boundaries, and uses truncated
backpropagation through time (TBPTT) during the K update epochs.

\subsection{World Models}

PlaNet \cite{hafner2019planet} introduced the Recurrent State-Space Model,
combining a deterministic GRU path with a stochastic latent and using
Cross Entropy Method planning in latent space. Dreamer
\cite{hafner2020dreamer} replaced planning with an actor--critic trained
on imagined latent rollouts, which requires a learned reward model.
Our implementation is deliberately more minimal: we keep only the
deterministic recurrent path, we supervise the belief state with both a
next-hidden-state prediction loss and a next-observation reconstruction
loss, and for imagination we rollout the transition head only (no reward
model) and update the policy with GAE bootstrapped from the critic on
zero-reward imagined trajectories. This simplification is a central object
of study in the paper.

\section{Related Work}

Memory-Gym \cite{pleines2023memorygym} benchmarks recurrent architectures
on memory-stressing grid worlds and finds that Recurrent PPO with
appropriate TBPTT length is a strong baseline. On \textsc{DoorKey}, Raileanu
and Rockt{\"a}schel's RIDE~\cite{raileanu2020ride} and related intrinsic
motivation methods demonstrate that exploration, rather than representation
learning, is the binding constraint as the room size grows. Our
cross-scale \textsc{DoorKey} results (both methods solve 5$\times$5;
neither solves 8$\times$8 reliably) are consistent with that literature:
auxiliary prediction losses that help when exploration is cheap do not
substitute for a dedicated exploration mechanism when it is not. The
Dreamer line of work
\cite{hafner2020dreamer, hafner2021dreamerv2, hafner2023dreamerv3} has
extended latent imagination to a wide range of domains; we view our study
as a diagnostic of what is lost when the stochastic latent, the KL balance,
and the learned reward predictor of that family are removed.

\section{Methods}

Both agents share a CNN image encoder, an action embedding, a GRU cell for
temporal integration, and an actor--critic head. They differ in whether
the belief state is supervised by an explicit world model and whether a
second update is performed on imagined rollouts.

\subsection{Shared Backbone}

The CNN follows the MiniGrid reference architecture from
\texttt{rl-starter-files}: three $2\times 2$ convolutions (16$\to$32$\to$64
channels) with a $2\times 2$ max-pool after the first conv, producing a
64-dimensional feature. A linear layer projects this to the 64-dimensional
observation embedding. Actions are embedded in a 16-dimensional space. The
GRU cell takes the concatenation of observation and previous-action
embeddings as input and produces the belief state $h_t$:
\begin{equation}
    h_t = \mathrm{GRU}\big(h_{t-1},\;[\text{CNN}(o_t);\; \text{Emb}(a_{t-1})]\big).
\end{equation}
The belief is reset to zero at episode boundaries. The actor and critic
are two-layer MLPs with Tanh activations and orthogonal initialization
(actor output gain $0.01$, critic output gain $1.0$).

\subsection{Recurrent PPO Baseline}

The baseline uses only the shared backbone. The objective is
\begin{equation}
    \mathcal{L} = \mathcal{L}_{\text{PPO}}
                  + c_v\,\mathcal{L}_V
                  - c_e\,H(\pi),
\end{equation}
with $\mathcal{L}_V = (V_\theta - R_t)^2$ and $H(\pi)$ the policy entropy.
The belief state is rolled out in the environment with gradient flow
through the recurrent cell, and TBPTT detaches the state every $\tau$ steps
to bound computation. Advantages are GAE-normalized.

\begin{algorithm}[!tbp]
\caption{Recurrent PPO --- one training update}
\label{alg:rppo}
\begin{algorithmic}[1]
\State Components: encoder $e_\omega$, belief $\mathrm{GRU}$, policy $\pi$,
value $V$.
\State \textbf{Rollout:} for $t = 1..T$ with $N$ parallel envs:
\State \quad $h_t \gets \mathrm{GRU}(h_{t-1}, [e_\omega(o_t)])$;
sample $a_t \sim \pi(\cdot \mid h_t)$; step envs.
\State \quad Store $(o_t, a_t, r_t, h_t, \log\pi(a_t \mid h_t))$ in replay buffer.
\State \quad Reset $h$ where episodes ended.
\State Compute bootstrap value $V(h_T)$; compute GAE advantages; normalize.
\State \textbf{PPO epoch} ($K$ times): replay trajectories, TBPTT every $\tau$ steps.
\State \quad Re-run $\mathrm{GRU}$ from stored $h_0$; get new $\log\pi_\theta(a_t \mid h_t)$,
$V_\phi(h_t)$.
\State \quad Compute
$\mathcal{L} = \mathcal{L}_{\text{CLIP}} + c_v \mathcal{L}_V - c_e H(\pi)$;
Adam + grad clip.
\State \quad $\mathcal{L}_{\text{CLIP}} = \mathbb{E}\!\left[\min\!\left(r_t \hat{A}_t,\,
\mathrm{clip}(r_t, 1{-}\varepsilon, 1{+}\varepsilon)\hat{A}_t\right)\right]$,
\hfill $r_t = \pi_\theta(a_t \mid h_t)\,/\,\pi_{\text{old}}(a_t \mid h_t)$
\end{algorithmic}
\end{algorithm}

\subsection{World-Model Agent}

The world-model agent adds two heads that share the belief state:
\begin{align}
    \hat h_{t+1} &= f_\phi(h_t,\, a_t) && \text{(transition)} \\
    \hat o_{t+1} &= g_\psi(h_t,\, a_t) && \text{(observation predictor)}
\end{align}
and the corresponding losses
$\mathcal{L}_{\text{trans}} = \lVert f_\phi(h_t, a_t) - h_{t+1}\rVert^2$
and
$\mathcal{L}_{\text{obs}} = \lVert g_\psi(h_t, a_t) - o_{t+1}\rVert^2$
(masked at episode boundaries). The joint objective is
\begin{equation}
    \mathcal{L} = \mathcal{L}_{\text{PPO}} + c_v\,\mathcal{L}_V
                  - c_e\,H(\pi)
                  + c_\text{t}\mathcal{L}_{\text{trans}}
                  + c_\text{o}\mathcal{L}_{\text{obs}}.
\end{equation}
Crucially, gradients from $\mathcal{L}_{\text{PPO}}$ flow back into the GRU
cell, so the belief state is shaped by the policy gradient as well as by
the prediction losses. An optional intrinsic reward
$r_\text{cur} = c_\text{cur}\lVert g_\psi(h_t, a_t) - o_{t+1}\rVert^2$
is added to the environment reward during rollout collection.

\paragraph{Imagination.}
After the first PPO update each epoch, we sample start beliefs from the
replayed trajectory, roll out $H_{\text{imag}}$ steps through the transition
model using the current policy (no environment calls, no reward model),
and perform a second Adam step with a REINFORCE-style policy term,
value error against GAE returns computed from zero imagined rewards plus
the critic's bootstrap at the horizon, and an entropy bonus.

\subsection{Training Loop}

Algorithm~\ref{alg:wm} summarizes one training iteration. Optimization uses
Adam with linear LR annealing, global gradient clipping at $0.5$, and
orthogonal weight initialization throughout. Rollouts are collected with
$N = 48$ vectorized environments over $T = 256$ steps, giving a batch of
$12{,}288$ transitions per update.

\begin{algorithm}[!tbp]
\caption{World Model Agent --- one training update}
\label{alg:wm}
\begin{algorithmic}[1]
\State Components: encoder $e_\omega$, belief $\mathrm{GRU}$, policy $\pi$,
value $V$, transition $f_\phi$, obs head $g_\psi$.
\State \textbf{Rollout:} for $t = 1..T$ with $N$ parallel envs:
\State \quad $h_t \gets \mathrm{GRU}(h_{t-1},[e_\omega(o_t);\mathrm{Emb}(a_{t-1})])$;
sample $a_t \sim \pi(\cdot\mid h_t)$; step envs.
\State \quad If $c_\text{cur}>0$: add
$c_\text{cur}\lVert g_\psi(h_t,a_t)-o_{t+1}\rVert^2$ to $r_t$.
\State \quad Reset $h$ where episodes ended.
\State Compute GAE advantages; normalize.
\State \textbf{PPO epoch} ($K$ times): replay trajectories, TBPTT every $\tau$ steps.
\State \quad Compute
$\mathcal{L} = \mathcal{L}_{\text{PPO}} + c_v\mathcal{L}_V - c_e H(\pi)
  + c_\text{t}\mathcal{L}_{\text{trans}} + c_\text{o}\mathcal{L}_{\text{obs}}$;
Adam + grad clip.
\State \textbf{Imagination} (after warmup): sample $h$ from buffer.
\State \quad For $j = 1..H_{\text{imag}}$: $a\sim\pi(\cdot\mid h)$,
$h \gets f_\phi(h, a)$.
\State \quad Second Adam step: REINFORCE policy term, value loss vs.\ GAE
on zero imagined rewards with $V(h)$ bootstrap, entropy bonus.
\end{algorithmic}
\end{algorithm}

\subsection{Hyperparameters}

Table~\ref{tab:hparams} lists the settings used for the main comparison.
These values are the result of the diagnostic work described in
Section~\ref{sec:diagnostics}; earlier configurations are reported there
as diagnosed failures.

\begin{table}[!tbp]
\centering
\footnotesize
\setlength{\tabcolsep}{4pt}
\begin{tabular}{lccc}
\toprule
 & \textsc{DoorKey-5x5} & \textsc{DoorKey-8x8} & \textsc{MemoryS11} \\
\midrule
Total steps & 3\,M & 10\,M & 10\,M \\
Parallel envs $N$ & 16 & 48 & 48 \\
Rollout length $T$ & 256 & 256 & 256 \\
Discount $\gamma$ & 0.99 & 0.99 & 0.99 \\
GAE $\lambda$ & 0.95 & 0.95 & 0.95 \\
Learning rate & $7\!\times\!10^{-4}$ & $3\!\times\!10^{-4}$ & $3\!\times\!10^{-4}$ \\
PPO epochs $K$ & 4 & 4 & 4 \\
Minibatches & 4 & 4 & 4 \\
Clip $\epsilon$ & 0.2 & 0.2 & 0.2 \\
Max grad norm & 0.5 & 0.5 & 0.5 \\
Max episode steps & 256 & 384 & 256 \\
Eval episodes & 30 & 30 & 100 \\
Entropy coef $c_e$ & 0.05 & 0.025 & 0.08 \\
Hidden dim & 128 & 256 & 256 \\
TBPTT length $\tau$ & 16 & 32 & 64 \\
Transition coef $c_\text{t}$ & 0.1 & 0.25 & 0.25 \\
Obs-pred coef $c_\text{o}$ & 0.1 & 0.25 & 0.25 \\
Curiosity coef $c_\text{cur}$ & 0.0 & 0.0 & 0.005 \\
Imagine horizon $H_\text{imag}$ & 0 & 0 & 0 \\
\bottomrule
\end{tabular}
\caption{Hyperparameters for the main comparison. \textsc{DoorKey-5x5} uses a dedicated sanity preset and a shorter budget because both methods converge early. The \textsc{DoorKey-8x8} column lists the preset we would use for a clean re-run; the actual \textsc{DoorKey-8x8} runs reported in Section~\ref{sec:experiments} come from an earlier exploratory phase that predated this preset (see the footnote in that paragraph for the exact per-run settings). The \textsc{MemoryS11} column lists the post-diagnostic preset used for the main MemoryS11 runs. Imagine horizon was set to 0 for final runs following the diagnostics in Section~\ref{sec:diagnostics}.}
\label{tab:hparams}
\end{table}

\section{Experiments}
\label{sec:experiments}

\subsection{Research Questions}

\begin{itemize}
    \item[\textbf{RQ1}] Does the world-model agent learn faster or reach
    higher return than Recurrent PPO on \textsc{MemoryS11} (recall-heavy
    POMDP)?
    \item[\textbf{RQ2}] Do the world-model auxiliary signals confer an
    advantage on \textsc{DoorKey}, and does that advantage survive as the
    exploration horizon grows from 5$\times$5 to 8$\times$8?
    \item[\textbf{RQ3}] Which components of the world-model pipeline
    (transition loss, observation reconstruction, curiosity, imagination)
    are individually responsible for the observed behaviour?
\end{itemize}

\subsection{Evaluation Protocol}

Every \texttt{eval\_interval}$\,=\,$10{,}000 environment steps we run
\texttt{eval\_episodes} stochastic on-policy evaluation episodes in a
separate vectorized env (30 for both \textsc{DoorKey} sizes, 100 for
\textsc{MemoryS11}) and log mean return and mean episode length. Actions
during evaluation are sampled from $\pi(\cdot\mid h_t)$ rather than taken
greedily, which keeps the reported return consistent with the on-policy
objective being optimized. Each configuration is run with multiple seeds;
reported curves show the seed mean with $\pm 1$ standard deviation shading.

\subsection{Main Comparison}

We organize the results by increasing difficulty: \textsc{DoorKey-5x5}
(both methods succeed), \textsc{MemoryS11} (both methods succeed,
world model learns more smoothly), and \textsc{DoorKey-8x8}
(neither method succeeds reliably).

\paragraph{\textsc{DoorKey-5x5}.}
Figure~\ref{fig:doorkey5} plots logged mean episode reward against
environment steps on the smaller \textsc{DoorKey} variant. Both curves
average seeds~0 and~42 (for the world model, seed~42 is linearly
interpolated onto seed~0's step grid before averaging). Both runs
approach near-optimal return under our settings; the averaged model-free
curve reaches mean eval return $\geq 0.95$ at approximately
$1.4\!\times\!10^5$ environment steps, whereas the averaged world-model
curve does so at approximately $3.0\!\times\!10^5$ steps, and both
converge to essentially the same asymptote (mean final return $0.965$
for Recurrent PPO, $0.961$ for the world-model agent). On this smallest
\textsc{DoorKey} variant the world-model agent is therefore \emph{not}
more sample-efficient than the recurrent baseline---it is in fact
slower to cross the $0.95$ line---but its curve is visibly smoother,
consistent with a regularizing effect of the auxiliary losses rather
than a sample-efficiency gain.

\begin{figure}[!tbp]
\centering
% PNG lives next to the repo under results/ (path is relative to this file in docs/).
\includegraphics[width=\columnwidth]{5x5.png}
\caption{\textsc{DoorKey-5x5-v0} training curves: mean episode reward vs.\
environment steps. Each line is the mean over seeds~0 and~42 for that
method.}
\label{fig:doorkey5}
\end{figure}

\paragraph{\textsc{MemoryS11}.}
Figure~\ref{fig:memory} shows the learning curves on \textsc{MemoryS11}.
Both methods converge to a mean eval return of approximately $0.47$
(Recurrent PPO $0.464\!\pm\!0.017$, world model $0.471\!\pm\!0.053$,
measured over the final $10\%$ of eval points, seed-averaged over
seeds~0 and~42), well short of the $\geq 0.9$ level that would indicate
reliable use of the remembered cue. We interpret this plateau as evidence
that the policy has learned to reach the correct half of the room more
often than chance (producing a return above the pure-random floor) but
has not learned a stable mapping from the early observation of the target
object to the correct later action. The two methods differ in the
\emph{shape} of learning rather than in its ceiling: the Recurrent PPO
curve reaches the plateau slightly earlier but with larger early
oscillations, while the world-model agent's curve climbs more slowly and
is smoother in the mid-training regime. Seed-to-seed variance at
convergence is comparable in magnitude for the two methods. The headline
finding, however, is that neither method reliably solves
\textsc{MemoryS11} under our configuration; the auxiliary world-model
objectives do not provide enough signal to push performance past the
apparent local optimum at $\approx 0.47$.

\begin{figure}[!tbp]
\centering
\includegraphics[width=\columnwidth]{memory.png}
\caption{\textsc{MemoryS11-v0} evaluation curves: mean eval return vs.\
environment steps. Each line averages seeds~0 and~42 for that method.}
\label{fig:memory}
\end{figure}

\paragraph{\textsc{DoorKey-8x8}.}
Figure~\ref{fig:doorkey} shows that neither method solves
\textsc{DoorKey-8x8} reliably within the training budget. The two
world-model seeds reported here were launched during the diagnostic
phase of the project and used different hyperparameter bundles from
each other and from the preset in Table~\ref{tab:hparams};
Table~\ref{tab:doorkey8_configs} lists the actual per-seed settings, which
are also preserved in the saved run configs. We therefore treat these
runs as an exploratory observation rather than a controlled comparison.
Mean return across seeds over the final $10\%$ of evaluations is
essentially $0$ for both methods ($0.000\!\pm\!0.000$ for Recurrent PPO,
$0.000\!\pm\!0.000$ for the world-model agent), and $0$ of $2$ seeds
reach goal-level return under either method within the training budget.
The stabilization effect we saw on 5$\times$5---where the world-model
agent produced a smoother learning curve at the same asymptote---does
not translate into solving 8$\times$8: when the room grows from
5$\times$5 to 8$\times$8, the expected number of random-policy steps
needed to execute the key-then-door sequence grows super-linearly, and
the auxiliary prediction losses---whose gradient magnitude does not
scale with task difficulty---do not substitute for a directed
exploration mechanism.
This is our most informative negative result, and we analyze it further
in Section~\ref{sec:discussion}.

\begin{figure}[!tbp]
\centering
\includegraphics[width=\columnwidth]{8x8.png}

\caption{Learning curves on \textsc{DoorKey-8x8-v0}. Neither method
reaches any success within the training budget: mean eval return stays
at $0$ across the two seeds per method that we report, consistent with
the well-known exploration bottleneck of sparse-reward
sequential-interaction tasks at this scale.}
\label{fig:doorkey}
\end{figure}

\begin{table}[!tbp]
\centering
\footnotesize
\setlength{\tabcolsep}{4pt}
\begin{tabular}{lcc}
\toprule
Hyperparameter       & Seed 0 & Seed 42 \\
\midrule
Learning rate        & $10^{-4}$ & $10^{-4}$ \\
Hidden dim           & 256       & 128       \\
TBPTT length $\tau$  & 32        & 16        \\
Rollout length $T$   & 128       & 128       \\
Imagine horizon      & 3         & 3         \\
Transition coef      & 0.25      & 0.1       \\
Obs-pred coef        & 0.25      & 0.1       \\
Entropy coef         & 0.025     & 0.05      \\
Curiosity coef       & 0         & 0.03      \\
\bottomrule
\end{tabular}
\caption{Actual per-seed hyperparameters used for the two
\textsc{DoorKey-8x8} world-model runs reported in this paper. These
differ both from each other and from the \textsc{DoorKey-8x8} column of
Table~\ref{tab:hparams}, which lists the preset we would use for a clean
re-run.}
\label{tab:doorkey8_configs}
\end{table}

\begin{table}[!tbp]
\centering
\footnotesize
\setlength{\tabcolsep}{3pt}
\begin{tabular}{lccc}
\toprule
Method & \textsc{DK-5x5} & \textsc{MemoryS11} & \textsc{DK-8x8} \\
\midrule
Recurrent PPO  & $0.965\!\pm\!0.00$ & $0.464\!\pm\!0.02$ & $0.000\!\pm\!0.00$ \\
World Model    & $0.961\!\pm\!0.00$ & $0.471\!\pm\!0.05$ & $0.000\!\pm\!0.00$ \\
\bottomrule
\end{tabular}
\caption{Mean final eval return (seed-averaged over seeds~0 and~42,
$\pm$ standard deviation of the seed-averaged curve over the last $10\%$
of eval points) at the end of training. \textsc{DK} abbreviates
\textsc{DoorKey}.}
\label{tab:final}
\end{table}

\subsection{Configuration Comparison on \textsc{MemoryS11}}
\label{sec:ablation}

A controlled single-factor ablation on \textsc{DoorKey-8x8} was not
feasible within our compute budget: each of our exploratory
\textsc{DoorKey-8x8} runs shipped with a different bundle of
hyperparameters (see the footnote in the \textsc{DoorKey-8x8} paragraph of
Section~\ref{sec:experiments}), which confounds component-level
attribution. We instead summarize the empirical effect of the two
hyperparameter changes that moved \textsc{MemoryS11} performance the most
during tuning, using pairs of runs from our training logs. These are
\emph{natural-configuration comparisons}, not randomized ablations: other
hyperparameters were not held perfectly fixed, so Table~\ref{tab:ablation}
should be read as ``the direction and approximate magnitude we observed
after a single-factor change,'' not ``the causal effect of flipping one
knob in isolation.''

\begin{table}[!tbp]
\centering
\footnotesize
\setlength{\tabcolsep}{4pt}
\begin{tabular}{p{0.26\columnwidth}p{0.42\columnwidth}p{0.20\columnwidth}}
\toprule
Change & Accompanying settings & Effect on return \\
\midrule
$H_\text{imag}{:}\,3{\to}0$ (turn off imagination)
  & also raised $c_\text{t}{=}c_\text{o}$ from $0.15/0.5$ to $0.25$
    and $c_e$ from $0.005$ to $0.05$--$0.08$
  & ${\approx}\,0.3{\to}0.5$ (stable) \\
$c_e{:}\,0.005{\to}0.05$ (raise entropy coef)
  & otherwise identical run family
  & ${\approx}\,0.15{\to}0.6$ (see Sec.~\ref{sec:diagnostics}) \\
\bottomrule
\end{tabular}
\caption{\textsc{MemoryS11} natural-configuration changes and their
observed effect on end-of-training mean evaluation return. Both rows
summarize multiple tuning runs rather than paired single-factor
ablations; the training logs backing each row are
\texttt{runs/world\_model\_MiniGrid-MemoryS11-v0\_*} and are preserved
with their \texttt{config.json} for inspection.}
\label{tab:ablation}
\end{table}

\subsection{Temporal Retention Probe on \textsc{MemoryS11}}
\label{sec:probe}

To distinguish \emph{belief-state failure} (the GRU forgets the cue) from
\emph{policy-head failure} (the GRU remembers the cue but the policy fails
to condition on it), we train a linear probe on the frozen belief state
of a trained \textsc{MemoryS11} checkpoint. This requires no additional
training runs.

\paragraph{Setup.}
We load the seed-0 \textsc{MemoryS11} world-model checkpoint
(hidden dim 256, post-diagnostic preset, mean eval return $\approx 0.47$),
roll it out for 400 fresh episodes under a seed band disjoint from
training (seed 20{,}000) and in-training evaluation (seed 10{,}000), and
record tuples $(h_t, y, t)$ where $y \in \{0,1\}$ is the identity of the
cue object (Ball vs.\ Key) placed in the starting room at grid position
$(1,4)$. We bin by step-within-episode and fit a logistic regression
probe from $h_t$ to $y$ per bin. Episodes are partitioned 70/30 at the
episode level (not at the timestep level) to prevent train--test
contamination. As a sanity check we run the identical pipeline with a
freshly initialized, untrained model.

\paragraph{Result.}
Table~\ref{tab:probe} reports probe test accuracy per bin together with
the per-bin majority-class baseline (unbalanced class frequencies in the
test split are a function of the trained policy's behaviour and so
differ across bins). In the cue-room phase the trained belief state lifts
probe accuracy about three points above majority, which we read as
``the cue is present in $h_t$ while the agent can still see it.'' From
the corridor onward, probe accuracy in the trained model drops to (or
below) majority: the cue is \emph{not} linearly decodable from $h_t$ at
the junction, where the decision is actually being made.

\begin{table}[!tbp]
\centering
\footnotesize
\setlength{\tabcolsep}{4pt}
\begin{tabular}{lcccc}
\toprule
Bin (step $t$) & $n_\text{test}$ & Maj. & \multicolumn{2}{c}{Probe test acc.} \\
\cmidrule(lr){4-5}
 & & base & Trained & Rand.\ init \\
\midrule
cue-room (0--4)       & 598  & 0.518 & 0.548 & 0.511 \\
corridor-early (5--14)  & 999  & 0.527 & 0.529 & 0.594 \\
corridor-mid (15--29)   & 659  & 0.619 & 0.470 & 0.571 \\
junction (30--59)      & 219  & 0.758 & 0.265 & 0.465 \\
late (60+)             & --   & --    & --    & 0.442 \\
\bottomrule
\end{tabular}
\caption{\textsc{MemoryS11} temporal retention probe: logistic regression
from $h_t$ to cue identity. ``Maj.\ base'' is the per-bin majority-class
accuracy on the test set; a probe with no information content would match
it. The trained checkpoint recovers the cue barely above majority only in
the cue-room phase; by the junction, where the cue must be used,
accuracy is at or below the majority baseline. The \emph{late} row is
populated only by the random-init baseline because the trained policy
typically terminates episodes before $t\!=\!60$.}
\label{tab:probe}
\end{table}

\paragraph{Interpretation.}
The probe is a direct test of whether the belief state carries the
information the junction decision would need; in our trained checkpoint
it does not. Combined with the fact that evaluation return plateaus
exactly at the level expected for a policy that cannot condition on the
cue, this pinpoints the failure as \emph{belief-state forgetting},
not policy-head miscalibration. It also complements the qualitative
finding in Section~\ref{sec:diagnostics} that detaching the GRU from the
policy gradient makes the task unsolvable: even with the GRU attached,
the current training regime does not shape the belief strongly enough
to preserve the cue across 30+ steps of corridor traversal.

\section{Diagnostic Experiments}
\label{sec:diagnostics}

Three design choices produced measurable regressions in early experiments
and became the most informative findings of the project. We report each
as a mini-experiment: hypothesis, quantitative result, mechanistic
diagnosis, and implication.

\subsection{BPTT Through the Transition Model in Imagination}

\paragraph{Hypothesis.} Letting policy and value gradients flow through the
transition model during imagination should make the dynamics task-relevant
rather than merely predictive, as in Dreamer.

\paragraph{Result.} Enabling imagination increased the transition loss by
roughly $9\times$ (\textbf{from $\approx 0.004$ to $\approx 0.036$}) once
the imagination phase engaged, and mean evaluation return regressed from
approximately $0.62$ to $0.30$ on \textsc{MemoryS11} within the same
training budget.

\paragraph{Diagnosis.} In our simplified setup the imagination rollout has
no learned reward model, so the only gradient signal reaching the policy
is the critic bootstrap $V(h_{t+H})$ evaluated on an imagined trajectory
with zero reward in its interior. This is a high-variance REINFORCE-style
estimator. When backpropagated through the transition network $f_\phi$,
this noisy signal competes with the clean one-step MSE prediction loss
and, empirically, dominates it: the transition model is pulled toward
whatever latent paths make the policy gradient cheaper, regardless of
physical consistency.

\paragraph{Implication.} Value-only imagination with BPTT through a
deterministic transition model is not a viable simplification of Dreamer.
Recovering the imagination benefit would require either (i) a learned
reward predictor---restoring the full Dreamer loop---or (ii) detaching the
transition network from the imagination policy gradient so that dynamics
prediction remains a clean auxiliary task. We adopt $H_\text{imag}\!=\!0$
in our final runs as the minimum-intervention fix.

\subsection{Entropy Regularization on Sparse-Reward Tasks}

\paragraph{Hypothesis.} A small entropy coefficient ($c_e\!=\!0.005$)
should let the policy commit confidently on \textsc{MemoryS11} once it
discovers the correct object--cue association.

\paragraph{Result.} Policy entropy plateaued at $\log 7 \approx 1.946$, i.e.,
the uniform distribution over actions, and evaluation return stuck near
$0.15$ for the entire training budget. Raising the coefficient to
$c_e\!=\!0.05$ dropped entropy to approximately $1.71$ and mean return
climbed past $0.6$.

\paragraph{Diagnosis.} With sparse binary reward, GAE advantages are
nearly zero for most samples; normalization does not fix this because it
rescales all advantages together. At $c_e\!=\!0.005$ the tiny policy
gradient is still dominated by the entropy bonus gradient, but with the
wrong sign: the policy is being pushed \emph{toward} uniform more strongly
than it is being pulled toward any particular action. Raising $c_e$ does
not ``encourage more exploration''---the policy is already maximally
exploratory---it changes the \emph{balance} so that signal from rare
rewarding trajectories is no longer drowned out.

\paragraph{Implication.} On sparse-reward POMDPs the
policy-gradient/entropy-gradient balance is a first-order hyperparameter,
not a fine-tuning detail. Sweeping $c_e$ across an order of magnitude is
probably more important than sweeping the learning rate or clip range.

\subsection{Detaching the GRU Hidden State from the Policy Gradient}

\paragraph{Hypothesis.} Detaching the belief state from the PPO loss
would keep PPO updates cheap and insulate the encoder from policy-gradient
noise, letting the reconstruction and transition losses shape
representations on their own.

\paragraph{Result.} \textsc{MemoryS11} was unsolvable under this
configuration across all seeds tried: mean return never moved above chance,
despite the reconstruction and transition losses converging normally.

\paragraph{Diagnosis.} The reconstruction and transition objectives are
\emph{content}-preserving: they reward the belief state for preserving
whatever is locally predictable about the next observation and the next
hidden state. Neither penalizes forgetting information that is not
immediately useful for prediction. In \textsc{MemoryS11} the critical
latent variable---the object seen at episode start---is locally irrelevant
for many time steps, so the reconstruction loss has no incentive to
maintain it in the belief state. Only the policy gradient, flowing back
through the GRU from the later action choice where the memory pays off,
can shape the belief to be \emph{task-useful}.

\paragraph{Implication.} Auxiliary world-model losses supplement, but
cannot replace, end-to-end gradient flow from the policy into the
recurrent cell. This is consistent with the behaviour of Dreamer and its
successors, all of which train the belief state with actor and critic
gradients in addition to reconstruction and KL terms.

\section{Discussion}
\label{sec:discussion}

\paragraph{When does the world model help?}
Our results across the three environments suggest that the benefit of a
learned world model tracks a single axis: how much the bottleneck is
\emph{representation} versus \emph{exploration}, but with an important
caveat that both axes can bite simultaneously. On \textsc{DoorKey-5x5},
where a random policy can stumble on the key-then-door sequence with
non-negligible probability, the auxiliary transition and observation
losses do not accelerate time-to-solve---the model-free baseline in fact
crosses the $0.95$ return line about twice as fast as the world-model
agent---but they do visibly smooth the learning curve and slightly
reduce seed-to-seed variance. Representation pressure is a real effect
on this task, it just shows up as a stability bonus rather than as a
sample-efficiency gain. On \textsc{MemoryS11}, the auxiliary reconstruction
objective does pressure the belief to encode visual features of the
recalled object while it is in view, but neither method pushes past a
mean return of $\approx 0.47$. Our temporal retention probe
(Section~\ref{sec:probe}) makes this concrete: the cue is barely above
majority-class decodable from $h_t$ while the agent is in the cue room,
and at or below majority by the junction. Representation-shaping from
the reconstruction loss is therefore not sufficient on its own to keep
the cue in the belief state through 30+ steps of corridor traversal,
which is exactly the horizon over which credit must be assigned to solve
the task. On
\textsc{DoorKey-8x8}, neither advantage matters: the problem is that the
agent rarely sees a successful trajectory at all. Auxiliary prediction
losses do not redistribute sparse rewards, and their gradient magnitude
does not grow with task difficulty. This pattern matches prior
observations in the MiniGrid literature that intrinsic motivation
specifically targeted at novel state visitation, rather than prediction
error per se, is what moves the needle on larger \textsc{DoorKey}
variants, and suggests that a similar directed-exploration component is
what \textsc{MemoryS11} needs to cross the $\approx 0.47 \to 0.9$ gap.

\paragraph{Auxiliary losses versus imagination.}
The two world-model ingredients played qualitatively different roles. The
auxiliary losses (transition and observation prediction) had a
regularizing effect: they did not raise the ceiling on \textsc{MemoryS11}
but they produced a smoother learning curve and lower per-step variance.
Imagination, at least in our value-only form, was actively harmful---the
BPTT-through-transition result in Section~\ref{sec:diagnostics} shows that
policy improvement on noisy imagined trajectories destabilizes the
dynamics model. A cleaner imagination design would separate the two
gradient paths: let the transition model be trained only by its one-step
prediction loss, and let the policy update on imagined rollouts without
propagating its gradient into $f_\phi$.

\paragraph{Limitations.}
Our implementation uses a purely deterministic belief with no stochastic
latent, no KL balance, and no learned reward model. These simplifications
were deliberate---we wanted to isolate the effect of the auxiliary
losses---but they also make our world-model agent a strictly weaker
version of Dreamer. The negative results on imagination should be read as
statements about this specific, minimalist design, not about
imagination-based RL in general. Compute constraints limited us to
2 seeds per configuration, and the \textsc{DoorKey} results
have high seed variance; a more rigorous study would use 5--10 seeds and
a larger training budget.

\section{Conclusion}

We compared Recurrent PPO against a deterministic world-model agent with
auxiliary prediction losses, optional curiosity, and Dyna-style
imagination on three partially observable MiniGrid tasks at two levels of
exploration difficulty. Both methods solve \textsc{DoorKey-5x5} to the
same asymptote (mean final return $\approx 0.96$), though on this task
the recurrent model-free baseline actually reaches the $0.95$ threshold
about twice as fast as the world-model agent; the world-model agent's
advantage on \textsc{DoorKey-5x5} is qualitative (smoother, lower-variance
curve), not a sample-efficiency win. On \textsc{MemoryS11} both methods
plateau at mean return $\approx 0.47$, short of reliable recall; the
world-model agent's curve is somewhat smoother in the mid-training
regime but the ceiling is the same. On \textsc{DoorKey-8x8} neither
method succeeds reliably---$0$ of $2$ seeds reach goal-level return
under either method. The world-model additions thus produce one
measurable qualitative effect---smoother learning---but do not accelerate
time-to-solve on any task and do not expand the set of tasks that are
solvable under our implementation. The most valuable
output of the project was a set of three diagnosed design trade-offs,
plus one post-hoc analysis of the belief state itself. The design
trade-offs are: value-only imagination with BPTT through a deterministic
transition model is unstable because it confounds dynamics learning with
a noisy policy gradient; under-regularized entropy on sparse-reward tasks
collapses policy learning rather than producing confident policies; and
detaching the belief state from the policy gradient makes memory tasks
unsolvable because auxiliary losses reward predictive utility, not task
utility. The analysis is a temporal retention probe on \textsc{MemoryS11}
that shows the cue identity is not linearly recoverable from the belief
state by the time the junction decision must be made, localizing the
$\approx 0.47$ plateau to belief-state forgetting rather than to a policy-head
failure.
Together these findings support a simple thesis: auxiliary world-model
objectives are useful as representation regularizers that pay off when
exploration is not the bottleneck, but they complement---they do not
substitute for---end-to-end policy-gradient training of the recurrent
belief state, and they cannot substitute for a dedicated exploration
mechanism when reward is sparse and the horizon is long.

\bibliography{refs}

\end{document}